% Generated draft: requires mathematical and source review.
\documentclass[11pt]{article}
\usepackage[margin=25mm]{geometry}
\usepackage{fontspec}
\setmainfont{Latin Modern Roman}
\usepackage{amsmath,amssymb,mathtools,unicode-math}
\setmathfont{Latin Modern Math}
\usepackage{enumitem,needspace,xurl,hyperref}
\setlength{\parindent}{0pt}
\setlength{\parskip}{5pt}
\setlength{\emergencystretch}{3em}
\newcommand{\N}{\mathbb{N}}
\newcommand{\Z}{\mathbb{Z}}
\newcommand{\Q}{\mathbb{Q}}
\newcommand{\R}{\mathbb{R}}
\newcommand{\C}{\mathbb{C}}
\newcommand{\F}{\mathbb{F}}
\newcommand{\A}{\mathbb{A}}
\newcommand{\PP}{\mathbb{P}}
\newcommand{\E}{\mathbb{E}}
\newcommand{\eps}{\varepsilon}
\newcommand{\dd}{\,\mathrm{d}}
\begin{document}
\section*{UnsolvedMath ID: 30000831}
\textbf{Source problem code:} catalogue entry
\subsection{Definitions and mathematical statement}
Let $(M,g)$ be a generic negatively curved Riemannian surface, as in the source. The word generic, the metric topology, and compactness/boundary hypotheses are not specified. To make one standard spectral convention explicit, suppose conditionally that $M$ is closed and use an orthonormal Laplace eigenbasis with \[-\Delta_g\phi_j=\lambda_j^2\phi_j,\qquad0\leq\lambda_0\leq\lambda_1\leq\cdots.\] Thus $\lambda$ is frequency, not the Laplace eigenvalue. Define \[N(\lambda)=\#\{j:\lambda_j\leq\lambda\},\quad N_{x,y}(\lambda)=\sum_{\lambda_j\leq\lambda}\phi_j(x)\overline{\phi_j(y)},\] \[R(\lambda)=N(\lambda)-\frac{\operatorname{Area}_g(M)}{4\pi}\lambda^2,\qquad R_x(\lambda)=N_{x,x}(\lambda)-\frac{1}{4\pi}\lambda^2.\] With these conventions, the source asks whether, for generic negatively curved surfaces, for every $\varepsilon>0$, \[R(\lambda)=O_{M,\varepsilon}(\lambda^\varepsilon),\qquad R_x(\lambda)=O_{M,x,\varepsilon}(\lambda^{1/2+\varepsilon}),\qquad N_{x,y}(\lambda)=O_{M,x,y,\varepsilon}(\lambda^{1/2+\varepsilon})\] as $\lambda\to\infty$, with the local estimate for every fixed $x\in M$ and the off-diagonal estimate for every fixed pair $x\ne y$. The last estimate is stated pointwise in the pair; no uniformity as $y\to x$ is inserted. If the source instead uses Laplace eigenvalues as its spectral parameter, the displayed definitions and exponents require corresponding renormalization, not silent reuse. Nor does the closed-surface illustration impose unverified compactness or boundary assumptions on the original fragment. The supplied 2026-08-21 triage labels the question open, reports logarithmic and averaged estimates and lower bounds, and says these do not establish the conjectured pointwise bounds for generic negatively curved surfaces.
\subsection{Short English statement}
Ask for subpolynomial global spectral remainder and frequency growth at most $\lambda^{1/2+\varepsilon}$ for each local remainder and fixed off-diagonal spectral kernel. The meaning of generic and the source's spectral and geometric conventions need confirmation.
\subsection{Sources}
\begin{itemize}\item Local UnsolvedMath registry, permanent ID 30000831, source problem catalogue entry, \emph{Spectral Remainders on Negatively Curved Surfaces}, proposed year 2007; no individual proposer supplied.\item Supplied citation: \emph{Analysis and Geometric Singularities} (2008), Oberwolfach Reports, DOI 10.4171/owr/2007/41, pp.\ 2361--2442: \url{https://doi.org/10.4171/owr/2007/41}. Supplied literature label: original Oberwolfach problem.\item Supplied literature label: current spectral remainder context, \url{https://arxiv.org/abs/2305.03549}.\end{itemize}
\end{document}
